% Generated draft: requires mathematical and source review.
\documentclass[11pt]{article}
\usepackage[margin=25mm]{geometry}
\usepackage{fontspec}
\setmainfont{Latin Modern Roman}
\usepackage{amsmath,amssymb,mathtools,unicode-math}
\setmathfont{Latin Modern Math}
\usepackage{enumitem,needspace,xurl,hyperref}
\setlength{\parindent}{0pt}
\setlength{\parskip}{5pt}
\setlength{\emergencystretch}{3em}
\newcommand{\N}{\mathbb{N}}
\newcommand{\Z}{\mathbb{Z}}
\newcommand{\Q}{\mathbb{Q}}
\newcommand{\R}{\mathbb{R}}
\newcommand{\C}{\mathbb{C}}
\newcommand{\F}{\mathbb{F}}
\newcommand{\A}{\mathbb{A}}
\newcommand{\PP}{\mathbb{P}}
\newcommand{\E}{\mathbb{E}}
\newcommand{\eps}{\varepsilon}
\newcommand{\dd}{\,\mathrm{d}}
\begin{document}
\section*{UnsolvedMath ID: 800021}
\textbf{Source problem code:} catalogue entry
\subsection{Definitions and mathematical statement}
A semigroup is a nonempty set with an associative binary operation. Let \(S\) be a commutative semigroup, written multiplicatively, and define
\[
\operatorname{End}(S)=\{f:S\to S: f(xy)=f(x)f(y)\text{ for all }x,y\in S\}.
\]
Equip this set with pointwise multiplication:
\[
(f\star g)(x)=f(x)g(x)\qquad(f,g\in\operatorname{End}(S),\ x\in S).
\]
Commutativity of \(S\) ensures that \(f\star g\) is again an endomorphism. Thus \((\operatorname{End}(S),\star)\) is a commutative semigroup. Endomorphisms here are semigroup homomorphisms; if \(S\) has an identity, no additional identity-preservation requirement is imposed.

Put
\[
\Q_{>0}=\{q\in\Q:q>0\},\qquad
\Q_{\geq0}=\{q\in\Q:q\geq0\}.
\]
\textbf{Problem (A. A. Borisov).} Describe, up to semigroup isomorphism, all commutative semigroups \(S\) satisfying each of the following conditions separately:
\begin{enumerate}
\item \((\operatorname{End}(S),\star)\cong(\Q_{>0},+)\);
\item \((\operatorname{End}(S),\star)\cong(\Q_{\geq0},+)\).
\end{enumerate}
Here each isomorphism is a bijection preserving the indicated semigroup operations. The operation on endomorphisms is explicitly pointwise multiplication. No finiteness, cancellation, or identity hypothesis on \(S\) is supplied.
\subsection{Short English statement}
Classify the commutative semigroups whose endomorphisms, multiplied pointwise, form a semigroup isomorphic to the positive rationals under addition, and separately classify those giving the nonnegative rationals under addition. The registry labels the problem open; this is a supplied status, not an independently verified assessment.
\subsection{Sources}
\begin{itemize}
\item Local UnsolvedMath registry, permanent ID \(800021\), source code \texttt{catalogue entry}, title \emph{The Lyapin's Notebook --- Problem 21}; supplied complete source record, upstream revision \texttt{d371ee917ac9647fbb06cd6fc529deec6f3dda82}.
\item Bershadsky, Kublanovsky, Mashevitzky, \emph{The Lyapin's Notebook: Unsolved Problems in Semigroup Theory} (2026), Problem 21. Supplied source URL: \url{https://arxiv.org/abs/2604.04763}. The record reports access on 2026-10-16 and extraction by \texttt{arxiv\_latex}; these provenance details are supplied and were not independently verified.
\item The problem statement attributes the problem to A. A. Borisov. The registry's \texttt{proposed\_by} field instead lists Bershadsky, Kublanovsky, Mashevitzky, with \texttt{proposed\_year} equal to 2026.
\end{itemize}
\end{document}
